\scanpage{130}\mspage{66}

\uline{VIII} \uline{Axiomatique : reformulation}.

Une structure \uline{triangulée} sur un ens $S$ est \inserted{la donnée d'}un ens
\[
\underset{\substack{\text{ens des}\\\text{faces}}}{F_0}\subset\unsure{\mathfrak{P}|\mathcal{P}|P}_3(S)\tag*{(1)}
\]
satisfaisant
\[
S=\bigcup_{f\in F}f\,,\tag*{(2)}
\]
et on lui demande ici de satisfaire les conditions

(3)\qquad (To 1)\qquad $F_0$ est l'ens des faces pour une structure de 2-polyèdre \deleted{\unsure{rég|rég.}} combinatoire \inserted{connexe} ayant $S$ comme ens des sommets, et comme ens des arêtes
\[
A_0=\{a\in\mathfrak{P}_2(S)\mid\exists f\in F_0,\ f\supset a\},
\]
et les relations d'incidences évidentes,

ce qui signifie que pour $\forall\,a\in\mathfrak{P}_2(S)$, s'il existe $f\ni F_0$ avec $f\supset a$, il en existe exactement 2, et que l'on a les conditions de connexité globale (deux \deleted{\unsure{sommets}}\inserted{faces} de $S$ sont reliables par une ``chaîne de faces'') et locale en tt $s\in S$ ($=$ chose, avec les faces $\ni s$). On suppose de plus que $F_0$ est déterminé par $\unsure{A_0|F_0}$ :
\[
\mathrm{To}\,2)\qquad F_0=\{f\in\mathfrak{P}_3(S))\mid\forall a\in\mathfrak{P}_2(f),\ \text{on a }a\in A_0\}.\tag*{(4)}
\]

Considérons maintenant les \deleted{deux}\deleted{\inserted{trois}}\inserted{deux} groupes

% archive sheet number: 129
